% Chapter opener (rewritten 2026-10-31 to the user's seven stages, Aïda-style single paragraph).
% Order of \input (no main file yet):
%   01a_inhouse_tracer      tracer on all BL scans (cited tool, not own)    sec:method-tracer
%   01b_medis_reference     Medis collapsed to full arteries; tracer vs Medis sec:method-medis
%   02_population_outcome   population, 17 segments, outcome, arteries      sec:method-population
%   02b_segment_split       cutting the traced arteries into segments       sec:method-split
%   02_models               prediction model, base terms (NOT WRITTEN)      sec:method-models
%   03_model_size           parameter budget and feature selection          sec:method-size,
%                           (Riley criteria: technical_background/03b)       sec:method-selection
%   11_cross_validation     cross-validation                                sec:method-cv
%   05_features to 09_      the chosen features in depth                    sec:method-features
%   10b_feature_validity    measurement properties                          sec:method-validity
% Retired, not \input: 10_medis_validation.tex (replaced by 01b_medis_reference.tex);
%   retired/04_feature_selection.tex (merged into 03_model_size.tex, 2026-11-06).
% Overview figure still to make (style guide: an overview figure opens the methods).
\providecommand{\clarify}[1]{\textcolor{red}{\textbf{Clarification needed:} #1}}
\chapter{Methodology}
\label{ch:methodology}

The methodology followed in this work is summarized in Figure~\ref{fig:method-overview} and consists of seven main stages:
(1) extraction of the coronary centerlines and lumen contours from all baseline scans with an in-house tracer;
(2) comparison of the extracted centerlines and lumen with Medis as a reference, as no ground truth was available;
(3) selection of the plaque-free baseline cohort and definition of the outcome from the visual analysis;
(4) division of the traced arteries into the segments of the visual analysis;
(5) the prediction model, the number of parameters the cohort can support, the choice of geometric features within it, and the cross-validation;
(6) definition of the geometric features and their expected measurement error;
and (7) assessment of the validity of these features against known groups and the variance explained by known risk factors.

\clarify{overview figure of the seven stages (label fig:method-overview).}
